\subsection{Recursion and free equational algebras}
\label{sec:annex-p10-recursion-free-algebras}

\begin{annexitem}{annex:p10-recursion-lifting}{thm:p10-recursion-lifting}{Transported recursion and the degenerate case}
\begin{proof}
Applying \(f^*\) or \(f_*\) to an endomap gives functors
between the two endomap categories.
The adjunction restricts to maps commuting with the endomaps:
transposing \(k f^*(s)=t k\) gives
\(k^\sharp s=f_*(t)k^\sharp\), and conversely.
Consequently, for every endomap \(t:B\to B\),
\[
 \begin{aligned}
 \operatorname{Hom}_{\operatorname{End}(\mathcal F)}
       ((f^*\omega,f^*s),(B,t))
 &\cong
 \operatorname{Hom}_{\operatorname{End}(\mathcal E)}
       ((\omega,s),(f_*B,f_*t))\\
 &\cong \operatorname{Hom}_{\mathcal E}(1,f_*B)\\
 &\cong \operatorname{Hom}_{\mathcal F}(1,B).
 \end{aligned}
\]
The middle bijection is the base recursion property, and the
last uses \(f^*1\cong1\).
The bijection is evaluation at the transported zero.
It is exactly the target recursion universal property.

To obtain free recursion with parameter object \(C\), let
\(h:C\to B\) and \(t:B\to B\).
Transpose \(h\) to a point of \(B^C\).
Postcomposition with \(t\) is an endomap of \(B^C\), so
recursion gives a unique map \(f^*\omega\to B^C\).
Uncurrying gives the unique map \(C\times f^*\omega\to B\)
with initial value \(h\) and commuting with successor and \(t\).
Thus \(C\times f^*\omega\) is free on \(C\).

It remains to check the extra nonequivalence clause.
In a nondegenerate topos the truth and falsity maps
\(1\to\Omega\) differ: they classify \(1\hookrightarrow1\)
and \(0\hookrightarrow1\), so equality would give \(0\cong1\).
The identity and the constant-true endomap of \(\Omega\)
therefore differ.
The identity underlying map from the first endomap object to
the second is not an endomap homomorphism.
Hence the forgetful functor is not full and cannot be an equivalence.

In a degenerate topos, \(0\cong1\).
Cartesian closedness makes product with any object preserve
the initial object, so every \(X\) is isomorphic to
\(X\times1\cong X\times0\cong0\).
There is only one morphism between any two objects.
Its endomap category is likewise equivalent to the terminal
category, and the forgetful functor is an equivalence.

The degenerate target cannot be excluded merely because a
geometric morphism was specified.
The constant functor from any topos to the terminal category
is left exact and left adjoint to the functor selecting its
terminal object. It is therefore the inverse-image functor
of a geometric morphism from the degenerate topos.
This is the stated exception to transporting the full
nonequivalence clause.
\end{proof}
\end{annexitem}

\begin{annexitem}{annex:p10-free-equational-algebras}{thm:p10-free-equational-algebras}{Terms quotients and their universal maps}
\begin{proof}
Fix the finite signature and let \(X\) be the generator object.

\emph{Finite words and parsing.}
The natural numbers object supplies finite ordinal objects
\([n]=\{k:k<n\}\) in a single family over \(n\in\omega\).
For any object \(Y\), dependent products form the family
of maps \([n]\to Y\); its total object is the object of
finite words over \(Y\). Empty words are included.
Concatenation, lengths, and finite scans are defined by
parameterized natural numbers recursion.
Thus this construction uses the already proved dependent
products, rather than an external countable coproduct.

Take the alphabet to be the disjoint sum of \(X\) and the
finite object of operation symbols, with their fixed arities.
Use postfix words. A variable pushes one item onto a stack.
An operation of arity \(k\) removes \(k\) items and pushes
one; a zero-ary operation simply pushes one.
A word is valid when no operation underflows and exactly
one item remains at the end.
These conditions are internal predicates obtained by scanning
the finite word and its natural-number stack depth.
They define a subobject \(T(X)\) of the word object.

The parsing is unambiguous. Scan instead with a stack of
already formed word blocks. A variable creates a singleton
block. An operation removes its ordered list of argument
blocks, concatenates them in order, and appends its own
symbol, yielding a single new block.
At each stage, concatenating the stack blocks gives the
input prefix already read. Underflow is exactly failure of
the depth condition. At a valid final stage there is one
block, namely the original word.
This determines its last constructor and all its immediate
argument words; each argument has strictly smaller length.
Conversely, concatenating valid argument words and appending
an operation gives precisely that constructor's valid word.
All operations on these stacks are finite-word operations.

\emph{The free term algebra.}
For an algebra \(B\) of the signature and a map \(h:X\to B\),
run the same valid-word scan with values in \(B\).
Adjoin a separate error state to the stack object and send underflow
to that state. This makes the scanning transition a total map,
even if \(B\) has no global elements.
Variables push their \(h\)-values, and an operation applies
the corresponding operation of \(B\) to the popped values.
Validity ensures that the required values are present;
the unique final value defines an arrow \(T(X)\to B\).
It commutes with each constructor by the parsing description.
Any constructor-preserving extension of \(h\) has that same
value, by induction on word length and the uniquely parsed
immediate arguments. Thus \(T(X)\) is the free algebra
before imposing the equations.

\emph{Imposing the equations.}
Each axiom has finitely many variables.
Substituting arbitrary elements of \(T(X)\) in those variables
gives an internal family of pairs of terms.
Intersect all equivalence relations on \(T(X)\) that contain
these pairs and are compatible with every operation.
The total relation is one such relation, and intersections
preserve all the stated conditions.
The least congruence therefore exists as a subobject of
\(T(X)\times T(X)\).
Its effective quotient exists by
\cref{thm:p10-lawvere-sheaves-effective-quotients}.

The operations descend because the congruence is compatible
with them and finite products of the quotient map are
epimorphisms. The quotient satisfies each equation:
precompose its two interpretation arrows with the
corresponding finite power of the quotient epimorphism.
They then agree by the imposed instances, so agree on
the quotient itself.
If \(B\) satisfies the equations, the kernel relation of
the term-evaluation map contains every imposed instance
and is an operation-compatible equivalence relation.
It contains the least congruence, so evaluation factors
uniquely through the quotient.
This proves the free-algebra universal property.
The induced maps on generators give its functoriality and
the free-forgetful adjunction.

Groups have a finite signature and finite equational
presentation using multiplication, inverse, and unit.
Unital rings have one using addition, additive inverse,
zero, multiplication, and one; commutativity adds one
equation. Thus all three named cases are covered.
No inequality \(0\neq1\) is imposed as an equational ring
axiom in this construction.

\emph{Change of topos.}
For these fixed finite signatures, \(f^*\) and \(f_*\)
transport operations and their equations using finite
products, including nullary operations.
Their adjunction restricts to algebra homomorphisms, by
transposing each operation-preservation equation.
For a free algebra on \(X\), that restricted adjunction gives
\[
 \operatorname{Hom}_{\mathrm{Alg}(\mathcal F)}
       (f^*\mathrm{Free}(X),B)
 \cong
 \operatorname{Hom}_{\mathcal F}(f^*X,UB),
\]
where \(U\) forgets the algebra structure.
It identifies \(f^*\mathrm{Free}(X)\) with the free algebra
on \(f^*X\), carrying the transported generator map.
The preceding recursion theorem ensures the target also
has the recursion object used in its term construction.
This proof does not assume preservation of power objects
or of the displayed impredicative description of the
least congruence.
\end{proof}
\end{annexitem}
